% Imporditud: tools/import_eksperimendid.py
% Allikas: Eesti füüsikaolümpiaadi eksperimendi ülesannete
%   kogu (Rael Kalda, 2023), ülesanne Ü58, lahendus L58.
\setAuthor{EFO žürii}
\setRound{piirkonnavoor}
\setYear{2012}
\setNumber{G E1}
\setDifficulty{2}
\setTopic{vedelike mehaanika}
\setVahendid{plastiliin, joonlaud, silindriline anum veega}
\setPunktid{10}
\setAllikas{Eksperimendi kogu Ü58, lahendus lk 58}
\setLahendusseis{toores}

\prob{Plastiliin}
Määrata plastiliini tihedus. Mõõtemääramatust ei ole vaja leida.
\textit{Märkus:} Sama ülesanne oli aastal 2001 gümnaasiumi piirkonnavoorus.

\hint

\solu
Kõigepealt tuleb sukeldada plastiliini veega täidetud silindrisse anumasse ja mõõta

joonlauaga seejuures toimuv veetaseme tõus $h_1$ (1 p). Seejärel tuleb valmistada plastiliinist laevuke ja mõõta veetaseme tõus $h_2$ siis, kui see laevuke ujub vee pinnal (2 p). Paneme tähele, et plastiliini ruumala $V_1$ = Sh1 (1 p), plastiliini mass m = $\rho$Sh2 (1 p) ja laevukese ujumisel välja surutava vee ruumala $V_2$ = Sh2 (1 p), mistõttu plastiliini mass m = $\rho$vSh2 (2 p). Nendest seostest saame avaldada $\rho$ = m/$V_1$ = $\rho$vh2/$h_1$ (1 p). Asendades siia mõõtmistulemused saame $\rho \approx$ 1,1 g/cm$^3$ kuni $\rho \approx$ 1,2 g/cm$^3$ (1 p) (tulemus oleneb plastiliini sordist).
\probend
